%%%%%%%%%%%%
%
% $Autor: Wings $
% $Datum: 2019-03-05 08:03:15Z $
% $Pfad: SetUp.tex $
% $Version: 4250 $
% !TeX spellcheck = en_GB/de_DE
% !TeX encoding = utf8
% !TeX root = manual 
% !TeX TXS-program:bibliography = txs:///biber
%
%%%%%%%%%%%%

\chapter{Inbetriebnahme}

Vor der ersten Nutzung muss ELMAR in Betrieb genommen werden. Nach Abschluss der folgenden Schritte ist das System betriebsbereit.

\begin{enumerate}
	
	\item Stellen Sie ELMAR auf eine ebene, nicht-brennbare und befestigte Oberfläche. 
	Achten Sie darauf, dass mindestens 30 cm Freiraum um ELMAR vorhanden sind.
	
	\item Überprüfen Sie, ob Arbeitsbereich und Antrieb frei sind.
	
	\item Überprüfen Sie, ob alle Kabel verbunden sind und der Netzschalter ausgeschaltet ist.
	
	\item Stecken Sie den Netzstecker ein.
	
	\item Schalten Sie den Netzschalter ein. ELMAR startet nun.
	
	\item Verbinden Sie Ihren Computer über ein USB-Kabel (USB-B) mit ELMAR.
	
	\item Starten Sie UGS.
	
	\item Wählen Sie in UGS den zugehörigen COM-Port aus und klicken Sie den "Connect or Disconnect" Button an, um eine Verbindung herzustellen.
	
	\item Überprüfen Sie, ob die Verbindung erfolgreich hergestellt wurde. Dies wird in der Steuerungssoftware durch eine Statusmeldung [Meldung] angezeigt.
	
	\item Führen Sie ein Homing durch.  
	
\end{enumerate}

Nach der erfolgreichen Referenzfahrt ist ELMAR betriebsbereit.